%%
%% Starter template for ASPLOS submissions.
%%
%% Source: https://www.asplos-conference.org/asplos2026/cfp/
%%
%% Preamble is as specified by the ASPLOS 2026 CFP.
%% Always verify against the current year's CFP before submitting.
%%

%% Use the base acmart.cls in sigplan proceedings format (10 pt default).
%% nonacm removes ACM-related text (copyright, CCS, etc.) for submission.
\documentclass[sigplan,nonacm]{acmart}

%% Enable page numbers.
\settopmatter{printfolios=true}

%% ============================================================
%% Packages
%% ============================================================
\usepackage{booktabs}    % professional tables
\usepackage{graphicx}    % \includegraphics
\usepackage{subcaption}  % subfigures
\usepackage{microtype}   % better text justification
\usepackage{xspace}      % smart spacing after macros
%% Placeholder tracking — copy utils/fillin.sty to this directory.
%% Use \fillin[desc], \fillcite[desc], \fillnum for unfinished values.
%% Add \printfillinlistifany before \end{document} during drafting.
%% Remove this line before submission.
\usepackage{fillin}

%% ============================================================
%% Custom macros
%% ============================================================
\newcommand{\sysname}{SysName\xspace}   % replace with your system's name
\newcommand{\eg}{e.g.,\xspace}
\newcommand{\ie}{i.e.,\xspace}
\newcommand{\etal}{\textit{et al.}\xspace}

%% ============================================================
\begin{document}

\title{\sysname: A [Noun Phrase Capturing Your Key Contribution]}

%% ASPLOS uses double-blind review. Omit author names for submission.
%% Restore for camera-ready.
\author{Author One}
\affiliation{\institution{Anonymous Institution}\country{}}
\author{Author Two}
\affiliation{\institution{Anonymous Institution}\country{}}

\begin{abstract}
%% 5-sentence formula:
%% 1. The problem (concrete, with a number if possible)
%% 2. Why existing systems fail to solve it
%% 3. Key insight + high-level design approach
%% 4. Evaluation setup: workload(s), baseline(s), scale
%% 5. Headline result

[1. Concrete problem statement.] [2. Why existing systems/hardware fail.]
[3. Key insight: we observe that X, enabling our design to do Y.] [4. We
implement \sysname in N~kLoC and evaluate it on [workload] against [baseline]
on [platform].] [5. \sysname achieves [Nx] improvement over [baseline] while
reducing [overhead] by [Y\%].]
\end{abstract}

\maketitle   % must come after the abstract (per ASPLOS CFP)

%% ============================================================
\section{Introduction}
%% ============================================================

%% Opening scenario: concrete situation where existing systems/hardware fail.
[Opening scenario with quantitative failure evidence.]

%% Fundamental challenge.
[Why existing approaches are structurally insufficient.]

%% Key insight (explicit sentence).
\textbf{Key insight.}
We observe that [one-sentence insight]. This enables [high-level design].

%% Brief system overview.
We design \sysname, [one-sentence description]. \S\ref{sec:design} describes it.

%% Contribution bullets (3–5).
This paper makes the following contributions:
\begin{itemize}
  \item We identify [observation] and show [evidence] (\S\ref{sec:motivation}).
  \item We design \sysname, achieving [goal] via [mechanism] (\S\ref{sec:design}).
  \item We implement \sysname in [N]k~LoC (\S\ref{sec:impl}).
  \item We show [Nx] improvement over [baseline] on [benchmark] (\S\ref{sec:eval}).
\end{itemize}

\S\ref{sec:motivation} motivates the work.
\S\ref{sec:design} describes the design.
\S\ref{sec:impl} covers implementation.
\S\ref{sec:eval} presents evaluation.
\S\ref{sec:related} discusses related work.
\S\ref{sec:conclusion} concludes.


%% ============================================================
\section{Motivation}
\label{sec:motivation}
%% ============================================================
%% ASPLOS papers often span hardware-software co-design.
%% Show the failure both at the application level AND at the
%% hardware/microarchitectural level where relevant.

\subsection{The Problem in Practice}

[Concrete scenario. Figure~\ref{fig:motivation} shows measured failure.]

\begin{figure}[t]
  \centering
  \includegraphics[width=\columnwidth]{figs/motivation.pdf}
  \caption{[Existing system/hardware] under [workload]: [metric] at [N\%]
           utilization. [Observation that motivates the work.]}
  \label{fig:motivation}
\end{figure}

\subsection{Why Existing Approaches Fall Short}

[Structural explanation. For architecture papers: include microarchitectural
analysis (e.g., cache miss rates, bandwidth bottlenecks, pipeline stalls).]

\subsection{Key Insight}

[Restate insight, connected to the evidence above.]


%% ============================================================
\section{Design}
\label{sec:design}
%% ============================================================

\sysname is designed to achieve:
\begin{enumerate}
  \item \textbf{Goal 1.} [Specific, measurable.]
  \item \textbf{Goal 2.} [Specific, measurable.]
\end{enumerate}
\sysname does \emph{not} target [non-goal].

\subsection{Overview}

[High-level description. Figure~\ref{fig:arch} shows the architecture.]

\begin{figure}[t]
  \centering
  \includegraphics[width=\columnwidth]{figs/architecture.pdf}
  \caption{\sysname architecture. [Component A] handles [X]; [Component B]
           manages [Y]. Shaded = new; unshaded = existing.}
  \label{fig:arch}
\end{figure}

\subsection{Hardware Interface}
%% If applicable: describe the hardware primitives or ISA extensions used.
[Hardware interface or assumptions. Remove this subsection if purely software.]

\subsection{Software Design}
[Software components and their interactions.]

\subsection{Correctness and Safety}
[Argue correctness properties relevant to your system.]


%% ============================================================
\section{Implementation}
\label{sec:impl}
%% ============================================================

We implement \sysname in [language], comprising [N]k~lines of code.
[Platform: OS, kernel version, hardware, FPGA, simulator, etc.]

\paragraph{[Key aspect 1.]} [Description.]
\paragraph{[Key aspect 2.]} [Description.]

%% If using simulation:
\paragraph{Simulation methodology.}
[Simulator name, version, validated against real hardware how?]


%% ============================================================
\section{Evaluation}
\label{sec:eval}
%% ============================================================

\subsection{Experimental Setup}

\paragraph{Platform.}
[CPU (model, GHz, cores, NUMA); Memory; Storage; Network; OS + kernel.
If simulation: simulator, modeled hardware configuration.]

\paragraph{Baselines.}
\begin{itemize}
  \item \textbf{[Baseline 1]} [version]: [why it is the right comparison.]
  \item \textbf{[Baseline 2]} [version]: [description.]
\end{itemize}

\paragraph{Workloads.}
[Benchmark suite and configuration. ASPLOS commonly uses: SPEC CPU 2017, PARSEC,
SPLASH-3, YCSB, TPC-C, Graph500, or domain-specific workloads.]

\subsection{Main Results}

\textbf{Claim:} \sysname achieves [Nx] improvement over [baseline] on [workload].

[Describe Figure~\ref{fig:main}. State what to observe and why.]

\begin{figure}[t]
  \centering
  \includegraphics[width=\columnwidth]{figs/main.pdf}
  \caption{[Metric] comparison on [workload]. \sysname achieves [Nx] over
           [baseline]. Error bars = one standard deviation over 5 runs.}
  \label{fig:main}
\end{figure}

\subsection{Sensitivity Analysis}
%% ASPLOS reviewers expect sensitivity to key parameters.
[How does performance vary with [parameter 1], [parameter 2]?]

\subsection{Overhead}
[Area overhead / power overhead / performance overhead on non-target workloads.]

\subsection{Microbenchmarks}
[Component-level benchmarks, each tied to a specific claim.]


%% ============================================================
\section{Related Work}
\label{sec:related}
%% ============================================================
%% Organize THEMATICALLY.

\paragraph{[Theme 1].} [...]
\paragraph{[Theme 2].} [...]


%% ============================================================
\section{Conclusion}
\label{sec:conclusion}
%% ============================================================

We present \sysname, [one-sentence summary]. Our key insight is that [restate].
\sysname achieves [headline result] over [baseline] on [workload].


%% Acknowledgments — remove or anonymize for submission
%\begin{acks}
%We thank [...]
%\end{acks}

%% Use the ACM bibliography style (per ASPLOS CFP)
\bibliographystyle{ACM-Reference-Format}
\bibliography{references}

\end{document}
